% c08 —— §8.4 单值性与扇（续）（p020–p021 / 原书 p.146–147）
% 承 c07 定理 8.2 之后的「ここから」；p021 末句「これらを順番に潰すことによって…」跨页，
% 整句归 c09。本块止于「……这些射影直线的法丛为 O(-1) 与 O(-5)。」句完。

% ===== 原书 p.146 =====

由此可得
\begin{align}
  & M_{\Delta_0} = \Pic Y_0 = E_8\perp E_8\perp
    \begin{pmatrix} 2 & 1 \\ 1 & -2 \end{pmatrix}, \tag{8.24}\\
  & M_{\Delta_0}^{\perp} = U\perp
    \begin{pmatrix} 2 & 1 \\ 1 & -2 \end{pmatrix}, \tag{8.25}\\
  & M_{\Delta_0} = \Pic\check{Y}_0 =
    \begin{pmatrix} 2 & 1 \\ 1 & -2 \end{pmatrix}, \tag{8.26}\\
  & M_{\Delta_0}^{\perp} = U\perp E_8\perp E_8\perp
    \begin{pmatrix} 2 & 1 \\ 1 & -2 \end{pmatrix}, \tag{8.27}
\end{align}
从而可知猜想 3.1 在此情形下成立。

% [Tutor Note: (8.24) 与 (8.26) 的左端同为 $M_{\Delta_0}$，分别指 $Y_0$ 与 $\check{Y}_0$ 的 Picard 格；(8.25)(8.27) 上标的 $\perp$ 按原图照录。]

$\check{X}$ 的 nef 锥由 $2D_1+D_4$ 与 $D_1$ 生成。这里，$D_1$ 与 $D_4$ 分别是与由
$v_1$ 与 $v_4$ 生成的一维锥对应的环面除子。把 $2D_1+D_4$ 与 $D_1$ 在 $\check{Y}$ 上的限制
分别记作 $E_1$ 与 $E_2$。对应的 $H^{2}_{\mathrm{amb}}(\check{Y},\mathbb{C}^{\times})$
在大体积极限附近的坐标 $(q_1,q_2)$，由
\begin{align}
  & q_1=\lambda, \tag{8.28}\\
  & q_2=\frac{\mu}{\lambda^{2}} \tag{8.29}
\end{align}
与二次栈的稠密环面 $\mathbb{L}^{\vee}_{\mathbb{C}^{\times}}\subset X_{\mathcal{F}(A_0)}$
的坐标 $(\lambda,\mu)$ 等同起来。

$\check{Y}$ 的代数格写作 $\mathcal{N}(\check{Y})\cong U\perp N$。这里，双曲平面
$U=\mathbb{Z}e\oplus\mathbb{Z}f$ 由 $e=[\mathcal{O}_p]$ 与
$f=-[\mathcal{O}_{\check{Y}}]-[\mathcal{O}_{E_1}]+3[\mathcal{O}_{E_2}]$ 生成。
正交补空间 $N\cong\left(\begin{smallmatrix}10&5\\5&2\end{smallmatrix}\right)$ 由
$e_1=[\mathcal{O}_{E_1}]$ 与 $e_2=[\mathcal{O}_{E_2}]$ 生成，并与 Picard 格
$\NS(\check{Y})$ 同构。$N$ 的正交群由两个元生成：
\begin{equation}
  O^{+}(N)=\langle g_1,g_2\rangle,\qquad
  g_1=\begin{pmatrix} 4 & 1 \\ -5 & -1 \end{pmatrix},\qquad
  g_2=\begin{pmatrix} 1 & 1 \\ 0 & -1 \end{pmatrix}. \tag{8.30}
\end{equation}

沿两条回路 $(q_1,q_2)\mapsto(e^{2\pi\sqrt{-1}}q_1,q_2)$ 与
$(q_1,q_2)\mapsto(q_1,e^{2\pi\sqrt{-1}}q_2)$ 的单值性 $T_1$ 与 $T_2$，由 (6.4)
分别给出为 $\mathcal{O}(-E_1)\otimes(-)$ 与 $\mathcal{O}(-E_2)\otimes(-)$。令
$\varphi_{e,\bullet}\colon N\hookrightarrow O(\mathcal{N}(\check{Y}))$ 为 (4.25) 给出的嵌入，则
\begin{align}
  & T_1=\varphi_{e,e_1}, \tag{8.31}\\
  & T_2=\varphi_{e,e_2}, \tag{8.32}
\end{align}
这容易由直接计算验证。

% [Tutor Note: 此处 $O(\mathcal{N}(\check{Y}))$ 是格 $\mathcal{N}(\check{Y})$ 的正交群，$O$ 按原图为斜体，与结构层 $\mathcal{O}$ 不同。]

如图 8.4 所示，把由 $e_1$ 与 $e_2$ 经 $O^{+}(N)$ 作用移动所得的向量生成的射线取作
一维锥，把以这些射线为一维锥的 $N_{\mathbb{R}}$ 中的扇记作 $\Sigma$。对应的环面簇
$X_{\Sigma}$ 是
% ===== 原书 p.147 =====
在 $(\mathbb{C}^{\times})^{2}$ 的边界上添加图 8.5 中由 $(-1)$ 直线与 $(-5)$ 直线交替排列
而成的无限长链而得到的。周期区域的环面紧化，通过把尖点附近替换为 $X_{\Sigma}$ 的边界
附近而得到。$X_{\Sigma}$ 上自然地作用着 $O^{+}(N)$，商空间 $X_{\Sigma}/O^{+}(N)$ 则先取
$X_{\Sigma}$ 关于由 $g_1$ 生成的无限循环群 $C_1\triangleleft O^{+}(N)$ 的商，再取关于由
$[g_2]$ 生成的 2 阶循环群 $C_2=O^{+}(N)/C_1$ 的商而得到。$g_1$ 在 $\Sigma$ 上的作用把一维锥
移到与其相隔两个位置的一维锥，因此把环面簇 $X_{\Sigma}$ 按此作用取商后，边界是把
$(-1)$ 曲线与 $(-5)$ 曲线如图 8.6 那样连成环状而得到的。商群 $C_2$ 在 $X_{\Sigma}/C_1$
上的作用把这个环沿水平线如图 8.7 那样翻转，于是由此得到的商成为图 8.8 中那样的加权
射影直线 $\mathbb{P}(1,2)$ 的链。这些射影直线的法丛为
$\mathcal{O}_{\mathbb{P}(1,2)}(-1)$ 与 $\mathcal{O}_{\mathbb{P}(1,2)}(-5)$。

% [Tutor Note: 原书「1 次元錐をその 2 つ隣の 1 次元錐に移す」的「2 つ隣」按字面译作「相隔两个位置」，即中间隔着一个一维锥。]

\begin{figure}[htbp]
  \centering
  \includegraphics[height=6.2cm]{fig8_4.png}
  \caption{扇 $\Sigma$（原书 p.147）}
\end{figure}

\begin{figure}[htbp]
  \centering
  \includegraphics[height=6.2cm]{fig8_5.png}
  \caption{环面簇 $X_\Sigma$（原书 p.147）}
\end{figure}

\begin{figure}[htbp]
  \centering
  \includegraphics[width=0.34\textwidth]{fig8_6.png}
  \caption{商簇 $X_\Sigma/C_1$（原书 p.147）}
\end{figure}

\begin{figure}[htbp]
  \centering
  \includegraphics[width=0.34\textwidth]{fig8_7.png}
  \caption{$C_2$ 的作用（原书 p.147）}
\end{figure}

\begin{figure}[htbp]
  \centering
  \includegraphics[width=0.34\textwidth]{fig8_8.png}
  \caption{商簇 $X_\Sigma/O^+(N)$（原书 p.147）}
\end{figure}

% CHUNK c08 DONE pages p020–p021
